\section{总结与反思}

\subsection{阶段成果}
项目将分散的教务、教学与校园生活入口放入同一应用。课表作为首页提供按学期及周次查询，学习页把跨教学班任务集中展示。本地待办构成可在离线条件下独立使用的场景，统一账号中心承担外部会话的操作入口。已核对的动态数据主要来自校方课表和有效成绩、学习通任务及798用水账户。现场可从离线课表和本地待办开始演示，具备有效会话时再进入在线查询。

\subsection{技术难点与解决思路}
教务登录循环使我认识到，用户授权和会话核验虽访问同一个门户，却不能共用界面恢复的触发条件。单测能够验证失效状态是否回写，只有覆盖安装后的实际启动和返回操作才能观察窗口是否再次弹出。798 的经历则说明，网络链路的验收应从真实服务结果倒推。申请短信、提交验证码、读取账户数据是三项不同的验证，某一环成功不应代替下一环的证据。将这些环节分开检查后，问题复现和修复验证可以使用相同的操作条件。

\subsection{AI 辅助开发与个人贡献}
我从江财学生换教室、查作业、记录个人安排和使用校园服务的场景出发，确定了五个模块的分工及统一账号入口。我确定课表优先回读、失效后主动重连的处理目标，也要求课程活跃度依据作业和通知判断，隐藏课程仍有机会在后续检查中恢复。开发中使用 Codex 辅助整理需求、阅读参考源码、编写及修改代码、排查故障、运行构建和定向测试，也辅助整理本报告及图表。我依据真实使用过程确认教务有效成绩的查询条件，亲自在授权页面完成必要登录，在真机复查798问题并反馈，发现单项调整仍无法登录后，锁定问题完成修复，并在新版上本人完成验证码验证，确认设备账号恢复。在校方返回空报表时，先回到真实页面核实查询条件；设备端持续登录失败时，提供可用参考包，把问题转为对照两个实际客户端。界面调整后则核对原事件绑定和前后截图。通过这些复核，代码修改可以对应到具体操作，假设、合成结果和真实服务结果也能在报告中分别说明。

\subsection{不足与后续计划}
目前需要优先补齐现场可复验的空白。教务应在校方开放真实非空考试安排后核对日期、地点和考试批次。学习通应继续观察真实新任务发布后的恢复行为，保留分项未读全时的提示。积分及设备启停涉及真实账户或扣费，验证须由用户本人发起。下一步继续开发洗浴原生能力，完善功能实现。
